% paper81-copilot-import-callable-advisors.tex — The Import and Callable Advisors:
%   Copilot-Guided Import Graph and Callable Surface Analysis
%   Paper 81 of the Judgment Geometry series.
% Compile: pdflatex paper81-copilot-import-callable-advisors
\documentclass[11pt]{article}
\usepackage{lmodern}
% jugeo-common.tex — shared preamble for all Judgment Geometry papers
% Include via: % jugeo-common.tex — shared preamble for all Judgment Geometry papers
% Include via: % jugeo-common.tex — shared preamble for all Judgment Geometry papers
% Include via: \input{jugeo-common}

% ── Packages ────────────────────────────────────────────────────────────
\usepackage{amsmath,amssymb,amsthm,mathtools}
\usepackage{tikz-cd}
\usepackage{listings}
\usepackage{xcolor}
\usepackage{xspace}
\usepackage{booktabs}
\usepackage{multirow}
\usepackage{enumitem}
\usepackage[expansion=false]{microtype}
\usepackage{hyperref}
\usepackage{cleveref}
% \usepackage{thmtools}  % disabled: not available in this TeX installation
\usepackage{float}
\usepackage{caption}
\usepackage{subcaption}
\usepackage{tabularx}
\usepackage{stmaryrd}       % \llbracket, \rrbracket
\usepackage{bbm}            % \mathbbm{1}

% ── Page geometry ───────────────────────────────────────────────────────
\usepackage[margin=1in]{geometry}

% ── Theorem environments ────────────────────────────────────────────────
\theoremstyle{plain}
\newtheorem{theorem}{Theorem}[section]
\newtheorem{lemma}[theorem]{Lemma}
\newtheorem{proposition}[theorem]{Proposition}
\newtheorem{corollary}[theorem]{Corollary}

\theoremstyle{definition}
\newtheorem{definition}[theorem]{Definition}
\newtheorem{example}[theorem]{Example}
\newtheorem{construction}[theorem]{Construction}

\theoremstyle{remark}
\newtheorem{remark}[theorem]{Remark}
\newtheorem{notation}[theorem]{Notation}

% ── Python listings style ──────────────────────────────────────────────
\definecolor{jugeoBlue}{HTML}{2563EB}
\definecolor{jugeoGreen}{HTML}{059669}
\definecolor{jugeoRed}{HTML}{DC2626}
\definecolor{jugeoPurple}{HTML}{7C3AED}
\definecolor{jugeoGray}{HTML}{6B7280}
\definecolor{jugeoBg}{HTML}{F8FAFC}

\lstdefinestyle{jugeo-python}{
  language=Python,
  basicstyle=\ttfamily\small,
  keywordstyle=\color{jugeoBlue}\bfseries,
  stringstyle=\color{jugeoGreen},
  commentstyle=\color{jugeoGray}\itshape,
  numberstyle=\tiny\color{jugeoGray},
  backgroundcolor=\color{jugeoBg},
  numbers=left,
  numbersep=6pt,
  frame=single,
  framerule=0.4pt,
  rulecolor=\color{jugeoGray!40},
  breaklines=true,
  breakatwhitespace=true,
  showstringspaces=false,
  tabsize=4,
  xleftmargin=1.5em,
  framexleftmargin=1.5em,
  morekeywords={dataclass,frozen,field,Enum,IntEnum,auto,Any,
                Mapping,Sequence,Optional,match,case},
  literate={→}{{$\to$}}1 {←}{{$\leftarrow$}}1
           {≤}{{$\leq$}}1 {≥}{{$\geq$}}1 {≠}{{$\neq$}}1
           {∈}{{$\in$}}1 {∀}{{$\forall$}}1 {∃}{{$\exists$}}1
           {⊕}{{$\oplus$}}1 {⊖}{{$\ominus$}}1,
  escapeinside={(*@}{@*)},
}
\lstset{style=jugeo-python}

\lstdefinestyle{jugeo-output}{
  basicstyle=\ttfamily\small,
  backgroundcolor=\color{jugeoBg},
  frame=single,
  framerule=0.4pt,
  rulecolor=\color{jugeoGray!40},
  breaklines=true,
  showstringspaces=false,
  xleftmargin=1.5em,
  framexleftmargin=0.5em,
  numbers=none,
}

% ── JuGeo-specific macros ──────────────────────────────────────────────

% Judgment 8-tuple
\newcommand{\J}{\mathcal{J}}
\newcommand{\Jud}[8]{(#1,\,#2,\,#3,\,#4,\,#5,\,#6,\,#7,\,#8)}
\newcommand{\JudShort}[1]{\J_{#1}}

% Components
\newcommand{\coord}{\mathsf{c}}            % coordinate
\newcommand{\prop}{\varphi}                 % proposition
\newcommand{\carrier}{\mathcal{A}}          % carrier type
\newcommand{\evid}{\mathcal{E}}             % evidence bundle
\newcommand{\oblig}{\mathcal{O}}            % obligations
\newcommand{\obstr}{\mathcal{B}}            % obstructions
\newcommand{\trust}{\mathcal{T}}            % trust annotation
\newcommand{\prov}{\Pi}                     % provenance

% Trust algebra
\newcommand{\Talg}{\mathfrak{T}}            % trust ordered algebra
\newcommand{\Eadm}{\mathcal{E}_{\mathrm{adm}}}
\newcommand{\tleq}{\preceq}                 % trust ordering
\newcommand{\tjoin}{\oplus}                 % conservative join
\newcommand{\tatten}{\ominus}               % attenuation
\newcommand{\tpromote}{\uparrow_\pi}        % promotion
\newcommand{\tdemote}{\downarrow_\chi}       % demotion

% Trust tiers
\newcommand{\Tcontra}{\ensuremath{\mathsf{CONTRADICTED}}}
\newcommand{\Tunver}{\ensuremath{\mathsf{UNVERIFIED}}}
\newcommand{\Tcopilot}{\ensuremath{\mathsf{COPILOT}}}
\newcommand{\Toracle}{\ensuremath{\mathsf{ORACLE}}}
\newcommand{\Truntime}{\ensuremath{\mathsf{RUNTIME}}}
\newcommand{\Tsolver}{\ensuremath{\mathsf{SOLVER}}}
\newcommand{\Tproof}{\ensuremath{\mathsf{PROOF}}}

% Site and geometry
\newcommand{\Site}{\mathcal{S}}             % semantic site
\newcommand{\Cov}{\mathrm{Cov}}            % covering family
\newcommand{\Ob}{\mathrm{Ob}}              % objects
\newcommand{\Mor}{\mathrm{Mor}}            % morphisms
\newcommand{\res}{\mathrm{res}}            % restriction
\newcommand{\Cech}{\check{C}}              % Čech complex
\newcommand{\Hcoh}[1]{H^{#1}}             % cohomology group

% Descent
\newcommand{\descent}{\mathrm{desc}}
\newcommand{\glue}{\mathrm{glue}}
\newcommand{\overlap}[2]{#1 \cap #2}

% Semantic moves
\newcommand{\VERIFY}{\textsc{Verify}}
\newcommand{\CONSTRUCT}{\textsc{Construct}}
\newcommand{\NORMALIZE}{\textsc{Normalize}}
\newcommand{\GLUE}{\textsc{Glue}}
\newcommand{\RESTRICT}{\textsc{Restrict}}
\newcommand{\TRANSPORT}{\textsc{Transport}}
\newcommand{\REPAIR}{\textsc{Repair}}
\newcommand{\PROMOTE}{\textsc{Promote}}
\newcommand{\CHALLENGE}{\textsc{Challenge}}
\newcommand{\ESCALATE}{\textsc{Escalate}}

% Fragments
\newcommand{\QFLIA}{\texttt{QF\_LIA}}
\newcommand{\QFLRA}{\texttt{QF\_LRA}}
\newcommand{\QFBV}{\texttt{QF\_BV}}
\newcommand{\QFUF}{\texttt{QF\_UF}}

% Misc
\newcommand{\Python}{\textsc{Python}}
\newcommand{\JuGeo}{\textsc{JuGeo}}
\newcommand{\LEAN}{\textsc{Lean}\,4}
\newcommand{\Fstar}{F$^{\star}$}
\newcommand{\Coq}{\textsc{Coq}}
\newcommand{\Dafny}{\textsc{Dafny}}

% Common shortcuts
\newcommand{\ie}{i.e.\@\xspace}
\newcommand{\eg}{e.g.\@\xspace}
\newcommand{\cf}{cf.\@\xspace}
\newcommand{\wrt}{w.r.t.\@\xspace}

% Evaluation shortcuts
\newcommand{\numcases}{300}
\newcommand{\numspec}{100}
\newcommand{\numeq}{100}
\newcommand{\numbug}{100}
\newcommand{\medtime}{6.5\,ms}
\newcommand{\totaltime}{1.949\,s}


% ── Packages ────────────────────────────────────────────────────────────
\usepackage{amsmath,amssymb,amsthm,mathtools}
\usepackage{tikz-cd}
\usepackage{listings}
\usepackage{xcolor}
\usepackage{xspace}
\usepackage{booktabs}
\usepackage{multirow}
\usepackage{enumitem}
\usepackage[expansion=false]{microtype}
\usepackage{hyperref}
\usepackage{cleveref}
% \usepackage{thmtools}  % disabled: not available in this TeX installation
\usepackage{float}
\usepackage{caption}
\usepackage{subcaption}
\usepackage{tabularx}
\usepackage{stmaryrd}       % \llbracket, \rrbracket
\usepackage{bbm}            % \mathbbm{1}

% ── Page geometry ───────────────────────────────────────────────────────
\usepackage[margin=1in]{geometry}

% ── Theorem environments ────────────────────────────────────────────────
\theoremstyle{plain}
\newtheorem{theorem}{Theorem}[section]
\newtheorem{lemma}[theorem]{Lemma}
\newtheorem{proposition}[theorem]{Proposition}
\newtheorem{corollary}[theorem]{Corollary}

\theoremstyle{definition}
\newtheorem{definition}[theorem]{Definition}
\newtheorem{example}[theorem]{Example}
\newtheorem{construction}[theorem]{Construction}

\theoremstyle{remark}
\newtheorem{remark}[theorem]{Remark}
\newtheorem{notation}[theorem]{Notation}

% ── Python listings style ──────────────────────────────────────────────
\definecolor{jugeoBlue}{HTML}{2563EB}
\definecolor{jugeoGreen}{HTML}{059669}
\definecolor{jugeoRed}{HTML}{DC2626}
\definecolor{jugeoPurple}{HTML}{7C3AED}
\definecolor{jugeoGray}{HTML}{6B7280}
\definecolor{jugeoBg}{HTML}{F8FAFC}

\lstdefinestyle{jugeo-python}{
  language=Python,
  basicstyle=\ttfamily\small,
  keywordstyle=\color{jugeoBlue}\bfseries,
  stringstyle=\color{jugeoGreen},
  commentstyle=\color{jugeoGray}\itshape,
  numberstyle=\tiny\color{jugeoGray},
  backgroundcolor=\color{jugeoBg},
  numbers=left,
  numbersep=6pt,
  frame=single,
  framerule=0.4pt,
  rulecolor=\color{jugeoGray!40},
  breaklines=true,
  breakatwhitespace=true,
  showstringspaces=false,
  tabsize=4,
  xleftmargin=1.5em,
  framexleftmargin=1.5em,
  morekeywords={dataclass,frozen,field,Enum,IntEnum,auto,Any,
                Mapping,Sequence,Optional,match,case},
  literate={→}{{$\to$}}1 {←}{{$\leftarrow$}}1
           {≤}{{$\leq$}}1 {≥}{{$\geq$}}1 {≠}{{$\neq$}}1
           {∈}{{$\in$}}1 {∀}{{$\forall$}}1 {∃}{{$\exists$}}1
           {⊕}{{$\oplus$}}1 {⊖}{{$\ominus$}}1,
  escapeinside={(*@}{@*)},
}
\lstset{style=jugeo-python}

\lstdefinestyle{jugeo-output}{
  basicstyle=\ttfamily\small,
  backgroundcolor=\color{jugeoBg},
  frame=single,
  framerule=0.4pt,
  rulecolor=\color{jugeoGray!40},
  breaklines=true,
  showstringspaces=false,
  xleftmargin=1.5em,
  framexleftmargin=0.5em,
  numbers=none,
}

% ── JuGeo-specific macros ──────────────────────────────────────────────

% Judgment 8-tuple
\newcommand{\J}{\mathcal{J}}
\newcommand{\Jud}[8]{(#1,\,#2,\,#3,\,#4,\,#5,\,#6,\,#7,\,#8)}
\newcommand{\JudShort}[1]{\J_{#1}}

% Components
\newcommand{\coord}{\mathsf{c}}            % coordinate
\newcommand{\prop}{\varphi}                 % proposition
\newcommand{\carrier}{\mathcal{A}}          % carrier type
\newcommand{\evid}{\mathcal{E}}             % evidence bundle
\newcommand{\oblig}{\mathcal{O}}            % obligations
\newcommand{\obstr}{\mathcal{B}}            % obstructions
\newcommand{\trust}{\mathcal{T}}            % trust annotation
\newcommand{\prov}{\Pi}                     % provenance

% Trust algebra
\newcommand{\Talg}{\mathfrak{T}}            % trust ordered algebra
\newcommand{\Eadm}{\mathcal{E}_{\mathrm{adm}}}
\newcommand{\tleq}{\preceq}                 % trust ordering
\newcommand{\tjoin}{\oplus}                 % conservative join
\newcommand{\tatten}{\ominus}               % attenuation
\newcommand{\tpromote}{\uparrow_\pi}        % promotion
\newcommand{\tdemote}{\downarrow_\chi}       % demotion

% Trust tiers
\newcommand{\Tcontra}{\ensuremath{\mathsf{CONTRADICTED}}}
\newcommand{\Tunver}{\ensuremath{\mathsf{UNVERIFIED}}}
\newcommand{\Tcopilot}{\ensuremath{\mathsf{COPILOT}}}
\newcommand{\Toracle}{\ensuremath{\mathsf{ORACLE}}}
\newcommand{\Truntime}{\ensuremath{\mathsf{RUNTIME}}}
\newcommand{\Tsolver}{\ensuremath{\mathsf{SOLVER}}}
\newcommand{\Tproof}{\ensuremath{\mathsf{PROOF}}}

% Site and geometry
\newcommand{\Site}{\mathcal{S}}             % semantic site
\newcommand{\Cov}{\mathrm{Cov}}            % covering family
\newcommand{\Ob}{\mathrm{Ob}}              % objects
\newcommand{\Mor}{\mathrm{Mor}}            % morphisms
\newcommand{\res}{\mathrm{res}}            % restriction
\newcommand{\Cech}{\check{C}}              % Čech complex
\newcommand{\Hcoh}[1]{H^{#1}}             % cohomology group

% Descent
\newcommand{\descent}{\mathrm{desc}}
\newcommand{\glue}{\mathrm{glue}}
\newcommand{\overlap}[2]{#1 \cap #2}

% Semantic moves
\newcommand{\VERIFY}{\textsc{Verify}}
\newcommand{\CONSTRUCT}{\textsc{Construct}}
\newcommand{\NORMALIZE}{\textsc{Normalize}}
\newcommand{\GLUE}{\textsc{Glue}}
\newcommand{\RESTRICT}{\textsc{Restrict}}
\newcommand{\TRANSPORT}{\textsc{Transport}}
\newcommand{\REPAIR}{\textsc{Repair}}
\newcommand{\PROMOTE}{\textsc{Promote}}
\newcommand{\CHALLENGE}{\textsc{Challenge}}
\newcommand{\ESCALATE}{\textsc{Escalate}}

% Fragments
\newcommand{\QFLIA}{\texttt{QF\_LIA}}
\newcommand{\QFLRA}{\texttt{QF\_LRA}}
\newcommand{\QFBV}{\texttt{QF\_BV}}
\newcommand{\QFUF}{\texttt{QF\_UF}}

% Misc
\newcommand{\Python}{\textsc{Python}}
\newcommand{\JuGeo}{\textsc{JuGeo}}
\newcommand{\LEAN}{\textsc{Lean}\,4}
\newcommand{\Fstar}{F$^{\star}$}
\newcommand{\Coq}{\textsc{Coq}}
\newcommand{\Dafny}{\textsc{Dafny}}

% Common shortcuts
\newcommand{\ie}{i.e.\@\xspace}
\newcommand{\eg}{e.g.\@\xspace}
\newcommand{\cf}{cf.\@\xspace}
\newcommand{\wrt}{w.r.t.\@\xspace}

% Evaluation shortcuts
\newcommand{\numcases}{300}
\newcommand{\numspec}{100}
\newcommand{\numeq}{100}
\newcommand{\numbug}{100}
\newcommand{\medtime}{6.5\,ms}
\newcommand{\totaltime}{1.949\,s}


% ── Packages ────────────────────────────────────────────────────────────
\usepackage{amsmath,amssymb,amsthm,mathtools}
\usepackage{tikz-cd}
\usepackage{listings}
\usepackage{xcolor}
\usepackage{xspace}
\usepackage{booktabs}
\usepackage{multirow}
\usepackage{enumitem}
\usepackage[expansion=false]{microtype}
\usepackage{hyperref}
\usepackage{cleveref}
% \usepackage{thmtools}  % disabled: not available in this TeX installation
\usepackage{float}
\usepackage{caption}
\usepackage{subcaption}
\usepackage{tabularx}
\usepackage{stmaryrd}       % \llbracket, \rrbracket
\usepackage{bbm}            % \mathbbm{1}

% ── Page geometry ───────────────────────────────────────────────────────
\usepackage[margin=1in]{geometry}

% ── Theorem environments ────────────────────────────────────────────────
\theoremstyle{plain}
\newtheorem{theorem}{Theorem}[section]
\newtheorem{lemma}[theorem]{Lemma}
\newtheorem{proposition}[theorem]{Proposition}
\newtheorem{corollary}[theorem]{Corollary}

\theoremstyle{definition}
\newtheorem{definition}[theorem]{Definition}
\newtheorem{example}[theorem]{Example}
\newtheorem{construction}[theorem]{Construction}

\theoremstyle{remark}
\newtheorem{remark}[theorem]{Remark}
\newtheorem{notation}[theorem]{Notation}

% ── Python listings style ──────────────────────────────────────────────
\definecolor{jugeoBlue}{HTML}{2563EB}
\definecolor{jugeoGreen}{HTML}{059669}
\definecolor{jugeoRed}{HTML}{DC2626}
\definecolor{jugeoPurple}{HTML}{7C3AED}
\definecolor{jugeoGray}{HTML}{6B7280}
\definecolor{jugeoBg}{HTML}{F8FAFC}

\lstdefinestyle{jugeo-python}{
  language=Python,
  basicstyle=\ttfamily\small,
  keywordstyle=\color{jugeoBlue}\bfseries,
  stringstyle=\color{jugeoGreen},
  commentstyle=\color{jugeoGray}\itshape,
  numberstyle=\tiny\color{jugeoGray},
  backgroundcolor=\color{jugeoBg},
  numbers=left,
  numbersep=6pt,
  frame=single,
  framerule=0.4pt,
  rulecolor=\color{jugeoGray!40},
  breaklines=true,
  breakatwhitespace=true,
  showstringspaces=false,
  tabsize=4,
  xleftmargin=1.5em,
  framexleftmargin=1.5em,
  morekeywords={dataclass,frozen,field,Enum,IntEnum,auto,Any,
                Mapping,Sequence,Optional,match,case},
  literate={→}{{$\to$}}1 {←}{{$\leftarrow$}}1
           {≤}{{$\leq$}}1 {≥}{{$\geq$}}1 {≠}{{$\neq$}}1
           {∈}{{$\in$}}1 {∀}{{$\forall$}}1 {∃}{{$\exists$}}1
           {⊕}{{$\oplus$}}1 {⊖}{{$\ominus$}}1,
  escapeinside={(*@}{@*)},
}
\lstset{style=jugeo-python}

\lstdefinestyle{jugeo-output}{
  basicstyle=\ttfamily\small,
  backgroundcolor=\color{jugeoBg},
  frame=single,
  framerule=0.4pt,
  rulecolor=\color{jugeoGray!40},
  breaklines=true,
  showstringspaces=false,
  xleftmargin=1.5em,
  framexleftmargin=0.5em,
  numbers=none,
}

% ── JuGeo-specific macros ──────────────────────────────────────────────

% Judgment 8-tuple
\newcommand{\J}{\mathcal{J}}
\newcommand{\Jud}[8]{(#1,\,#2,\,#3,\,#4,\,#5,\,#6,\,#7,\,#8)}
\newcommand{\JudShort}[1]{\J_{#1}}

% Components
\newcommand{\coord}{\mathsf{c}}            % coordinate
\newcommand{\prop}{\varphi}                 % proposition
\newcommand{\carrier}{\mathcal{A}}          % carrier type
\newcommand{\evid}{\mathcal{E}}             % evidence bundle
\newcommand{\oblig}{\mathcal{O}}            % obligations
\newcommand{\obstr}{\mathcal{B}}            % obstructions
\newcommand{\trust}{\mathcal{T}}            % trust annotation
\newcommand{\prov}{\Pi}                     % provenance

% Trust algebra
\newcommand{\Talg}{\mathfrak{T}}            % trust ordered algebra
\newcommand{\Eadm}{\mathcal{E}_{\mathrm{adm}}}
\newcommand{\tleq}{\preceq}                 % trust ordering
\newcommand{\tjoin}{\oplus}                 % conservative join
\newcommand{\tatten}{\ominus}               % attenuation
\newcommand{\tpromote}{\uparrow_\pi}        % promotion
\newcommand{\tdemote}{\downarrow_\chi}       % demotion

% Trust tiers
\newcommand{\Tcontra}{\ensuremath{\mathsf{CONTRADICTED}}}
\newcommand{\Tunver}{\ensuremath{\mathsf{UNVERIFIED}}}
\newcommand{\Tcopilot}{\ensuremath{\mathsf{COPILOT}}}
\newcommand{\Toracle}{\ensuremath{\mathsf{ORACLE}}}
\newcommand{\Truntime}{\ensuremath{\mathsf{RUNTIME}}}
\newcommand{\Tsolver}{\ensuremath{\mathsf{SOLVER}}}
\newcommand{\Tproof}{\ensuremath{\mathsf{PROOF}}}

% Site and geometry
\newcommand{\Site}{\mathcal{S}}             % semantic site
\newcommand{\Cov}{\mathrm{Cov}}            % covering family
\newcommand{\Ob}{\mathrm{Ob}}              % objects
\newcommand{\Mor}{\mathrm{Mor}}            % morphisms
\newcommand{\res}{\mathrm{res}}            % restriction
\newcommand{\Cech}{\check{C}}              % Čech complex
\newcommand{\Hcoh}[1]{H^{#1}}             % cohomology group

% Descent
\newcommand{\descent}{\mathrm{desc}}
\newcommand{\glue}{\mathrm{glue}}
\newcommand{\overlap}[2]{#1 \cap #2}

% Semantic moves
\newcommand{\VERIFY}{\textsc{Verify}}
\newcommand{\CONSTRUCT}{\textsc{Construct}}
\newcommand{\NORMALIZE}{\textsc{Normalize}}
\newcommand{\GLUE}{\textsc{Glue}}
\newcommand{\RESTRICT}{\textsc{Restrict}}
\newcommand{\TRANSPORT}{\textsc{Transport}}
\newcommand{\REPAIR}{\textsc{Repair}}
\newcommand{\PROMOTE}{\textsc{Promote}}
\newcommand{\CHALLENGE}{\textsc{Challenge}}
\newcommand{\ESCALATE}{\textsc{Escalate}}

% Fragments
\newcommand{\QFLIA}{\texttt{QF\_LIA}}
\newcommand{\QFLRA}{\texttt{QF\_LRA}}
\newcommand{\QFBV}{\texttt{QF\_BV}}
\newcommand{\QFUF}{\texttt{QF\_UF}}

% Misc
\newcommand{\Python}{\textsc{Python}}
\newcommand{\JuGeo}{\textsc{JuGeo}}
\newcommand{\LEAN}{\textsc{Lean}\,4}
\newcommand{\Fstar}{F$^{\star}$}
\newcommand{\Coq}{\textsc{Coq}}
\newcommand{\Dafny}{\textsc{Dafny}}

% Common shortcuts
\newcommand{\ie}{i.e.\@\xspace}
\newcommand{\eg}{e.g.\@\xspace}
\newcommand{\cf}{cf.\@\xspace}
\newcommand{\wrt}{w.r.t.\@\xspace}

% Evaluation shortcuts
\newcommand{\numcases}{300}
\newcommand{\numspec}{100}
\newcommand{\numeq}{100}
\newcommand{\numbug}{100}
\newcommand{\medtime}{6.5\,ms}
\newcommand{\totaltime}{1.949\,s}


% ── Additional packages ────────────────────────────────────────────
\usepackage{array}
\usepackage{tikz}
\usetikzlibrary{arrows.meta,positioning,calc,fit,backgrounds}

% ── Paper-specific macros ──────────────────────────────────────────
\newcommand{\ImportAdv}{\mathsf{CopilotImportAdvisor}}
\newcommand{\CallAdv}{\mathsf{CopilotCallableAdvisor}}
\newcommand{\ImportGraph}{\mathsf{ImportGraph}}
\newcommand{\CallSurf}{\mathsf{CallableSurface}}
\newcommand{\CycleDetect}{\mathrm{detectCycles}}
\newcommand{\StarImport}{\mathrm{starImport}}
\newcommand{\BindingErr}{\mathrm{bindingError}}
\newcommand{\DescAnalysis}{\mathrm{descriptorAnalysis}}
\newcommand{\AdvSound}{\mathrm{sound}}
\newcommand{\CopChan}{\mathsf{CopilotChannel}}
\newcommand{\CopBridge}{\mathsf{CopilotAPIBridge}}
\newcommand{\TrustCeil}{\mathrm{ceiling}}
\newcommand{\ImportCycle}{\mathsf{ImportCycle}}
\newcommand{\CallTarget}{\mathsf{CallTarget}}
\newcommand{\AdvCompose}{\otimes}

\title{\textbf{The Import and Callable Advisors:\\
Copilot-Guided Import Graph and Callable Surface Analysis}}

\author{Halley Young}

\date{Paper~81 --- Judgment Geometry Series}

\begin{document}
\maketitle

% ─────────────────────────────────────────────────────────────────────
\begin{abstract}
Python's dynamic import system and flexible calling conventions
introduce verification challenges that static analyses alone
cannot fully address.  Circular imports may cause runtime failures,
star imports pollute namespaces, and descriptor-based calling
conventions can lead to subtle binding errors.  We introduce two
Copilot-guided advisors within the \JuGeo{} framework:
(i)~the $\ImportAdv$, which analyzes import graphs for cycles,
star imports, and dependency anomalies; and
(ii)~the $\CallAdv$, which analyzes callable surfaces for
descriptor resolution errors, incorrect binding, and signature
mismatches.

Both advisors query the $\CopChan$ for LLM-generated analysis
hints and produce advice at the $\Tcopilot$ trust tier.  We
formalize the import graph as a directed graph with cycle detection
semantics and the callable surface as a resolution lattice.  We
prove five formal properties: import advisory soundness, callable
advisory soundness, cycle detection completeness, binding error
monotonicity, and composition trust preservation.  All proofs are
presented in full and formalized in \LEAN.
\end{abstract}

% ─────────────────────────────────────────────────────────────────────
\section{Introduction}
\label{sec:intro}

Python's import system is one of the most flexible---and most
error-prone---module systems in widespread use.  Circular imports,
conditional imports, star imports, and dynamic import expressions
create a complex dependency landscape that defies simple static
analysis.  Similarly, Python's callable surface---the set of
objects that can be invoked as functions---is complicated by
descriptors, metaclasses, \texttt{\_\_call\_\_} methods, and
dynamic dispatch.

The \JuGeo{} framework addresses these challenges through
domain-specific advisors that combine LLM heuristics with formal
verification.  The $\ImportAdv$ focuses on the import graph,
while the $\CallAdv$ focuses on callable surfaces.  Both advisors
follow the advisor pattern established in Paper~90: they generate
advice at the $\Tcopilot$ trust tier and require independent
validation before influencing the judgment state.

\paragraph{Contributions.}
\begin{itemize}[leftmargin=2em]
\item We formalize the Python import graph as a directed graph and
  define cycle detection, star import analysis, and dependency
  anomaly detection (\S\ref{sec:import}).

\item We formalize the callable surface as a resolution lattice
  and define descriptor analysis, binding error detection, and
  signature validation (\S\ref{sec:callable}).

\item We prove soundness, completeness, and trust preservation
  for both advisors (\S\ref{sec:formal}).

\item We demonstrate integration with the \JuGeo{} API and
  evidence channel infrastructure (\S\ref{sec:integration}).
\end{itemize}

\paragraph{Outline.}
Section~\ref{sec:pattern} reviews the advisor pattern.
Section~\ref{sec:import} presents the import advisor.
Section~\ref{sec:callable} presents the callable advisor.
Section~\ref{sec:integration} covers API integration.
Section~\ref{sec:formal} proves the main theorems.
Section~\ref{sec:discussion} discusses implications.
Section~\ref{sec:related} surveys related work.
Section~\ref{sec:conclusion} concludes.

% ─────────────────────────────────────────────────────────────────────
\section{The Advisor Pattern}
\label{sec:pattern}

We briefly recall the advisor pattern from Paper~90.

\begin{definition}[Advisor]
\label{def:advisor}
An advisor $A$ accepts a verification context, queries the
$\CopChan$, and produces advice items at trust $\Tcopilot$.
Each advice item $a = (d, s, \gamma, \Tcopilot)$ carries a
domain $d$, suggestion $s$, confidence $\gamma$, and fixed trust.
\end{definition}

\begin{definition}[Advisory validation]
\label{def:validation}
An advice item $a$ is \emph{validated} if the solver or runtime
channel independently confirms the issue or suggestion.
Validated advice may be promoted to higher trust.
\end{definition}

% ─────────────────────────────────────────────────────────────────────
\section{The Copilot Import Advisor}
\label{sec:import}

\subsection{Import Graph Formalization}
\label{sec:import:graph}

\begin{definition}[Import graph]
\label{def:import-graph}
The \emph{import graph} $G = (V, E)$ of a Python project has:
\begin{itemize}[leftmargin=2em]
\item Vertices $V$: the set of Python modules (files).
\item Edges $E$: directed edges $(m_1, m_2)$ indicating that
  module $m_1$ imports from module $m_2$.
\end{itemize}
An edge is \emph{unconditional} if the import is at module level,
and \emph{conditional} if it is inside a function or guarded by
a condition.
\end{definition}

\begin{definition}[Import cycle]
\label{def:import-cycle}
An \emph{import cycle} is a sequence of modules
$m_1 \to m_2 \to \cdots \to m_k \to m_1$ where each edge
is an unconditional import.  A cycle is \emph{benign} if
all modules involved have already been fully initialized by
the time the cycle is traversed, and \emph{hazardous} otherwise.
\end{definition}

\begin{definition}[Star import]
\label{def:star-import}
A \emph{star import} is an import of the form
$\texttt{from M import *}$, which imports all public names
from module $M$ into the current namespace.  A star import is
\emph{safe} if $M$ defines $\texttt{\_\_all\_\_}$ and all
exported names are non-conflicting.
\end{definition}

The following listing shows how the $\ImportAdv$ uses the
$\CopChan$ for import analysis.

\begin{lstlisting}[style=jugeo-python, caption={Import advisor querying the CopilotChannel.}, label=lst:import-query]
from jugeo.evidence.channels import (
    CopilotChannel,
    ChannelConfiguration,
    EvidenceChannel,
    ChannelJurisdiction,
    EvidenceRequest,
    EvidenceResponse,
    ChannelMonitor,
)
from jugeo.interfaces.api import (
    JuGeoAPI,
    APIRequest,
    APIResponse,
    OperationKind,
    RequestStatus,
    CopilotAPIBridge,
    APIEventLog,
)

# Configure for import analysis
config = ChannelConfiguration(
    channel=EvidenceChannel.COPILOT,
    timeout_ms=10000,
    max_retries=2,
    batch_size=1,
    rate_limit=8.0,
    trust_ceiling="proposal",
    jurisdiction=ChannelJurisdiction.for_channel(
        EvidenceChannel.COPILOT
    ),
    is_enabled=True,
    priority=40,
)
copilot = CopilotChannel(config=config, model_name="copilot-default")

# Query for import cycle detection
import_prompt = (
    "Analyze the import structure:\n"
    "  # module_a.py\n"
    "  from module_b import helper\n"
    "  # module_b.py\n"
    "  from module_c import util\n"
    "  # module_c.py\n"
    "  from module_a import config  # potential cycle\n"
    "Identify: (1) import cycles, (2) star imports,\n"
    "(3) dependency anomalies."
)
result = copilot.query_llm(prompt=import_prompt, timeout_ms=8000)

# Parse the response
parsed = copilot.parse_response()
copilot.apply_trust_ceiling()

# Build evidence request for import analysis
request = EvidenceRequest(
    request_id="import-advisory-001",
    coordinate="project.imports.cycle",
    proposition="no hazardous circular imports exist",
    required_kind="PROPOSAL",
    proposition_kind="import_analysis",
    preferred_channel=EvidenceChannel.COPILOT,
    fallback_channels=(EvidenceChannel.RUNTIME,),
    deadline_ms=10000.0,
    budget=2.0,
    metadata={"advisor": "import", "check": "cycles"},
)

monitor = ChannelMonitor()
monitor.record_request(request)
response = copilot.handle_request(request)
monitor.record_response(response)

print(f"Import advisory trust: {response.trust_level}")
print(f"Is partial: {response.is_partial}")
print(f"Residuals: {response.residuals}")
print(f"Success rate: {monitor.success_rate()}")

# API integration for import inspection
event_log = APIEventLog()
bridge = CopilotAPIBridge(
    channel=copilot,
    event_log=event_log,
    require_corroboration=True,
)
api = JuGeoAPI(event_log=event_log, copilot_bridge=bridge)

inspect_req = APIRequest(
    request_id="import-inspect-001",
    operation=OperationKind.INSPECT,
    coordinate="project.imports",
    proposition="import graph is acyclic",
    budget=8000,
    deadline=15.0,
    metadata={"advisor": "import"},
)
inspect_resp = api.inspect(inspect_req)
print(f"Inspection: {inspect_resp.status}")
\end{lstlisting}

\subsection{Import Advisory Analysis}
\label{sec:import:analysis}

The $\ImportAdv$ performs three types of analysis:

\begin{enumerate}
\item \textbf{Cycle detection.}  The LLM identifies potential
  import cycles and classifies them as benign or hazardous.

\item \textbf{Star import analysis.}  The LLM identifies star
  imports and checks whether they are safe (module defines
  $\texttt{\_\_all\_\_}$) or potentially problematic.

\item \textbf{Dependency anomaly detection.}  The LLM identifies
  unusual import patterns: conditional imports that may fail,
  imports of deprecated modules, and import-time side effects.
\end{enumerate}

\begin{remark}[Import analysis depth]
\label{rem:import-depth}
The $\ImportAdv$ can operate at varying depths: shallow analysis
considers only direct imports, while deep analysis follows the
transitive import closure.  The depth is controlled by the
budget parameter in the evidence request.
\end{remark}

% ─────────────────────────────────────────────────────────────────────
\section{The Copilot Callable Advisor}
\label{sec:callable}

\subsection{Callable Surface Formalization}
\label{sec:callable:surface}

\begin{definition}[Callable surface]
\label{def:callable-surface}
The \emph{callable surface} of a Python program is the set of
all objects that can be invoked.  This includes:
\begin{itemize}[leftmargin=2em]
\item Functions (defined with \texttt{def}).
\item Lambda expressions.
\item Class instances with \texttt{\_\_call\_\_} methods.
\item Bound and unbound methods.
\item Descriptors (\texttt{\_\_get\_\_}, \texttt{\_\_set\_\_}).
\item Metaclass-instantiated callables.
\end{itemize}
\end{definition}

\begin{definition}[Callable resolution]
\label{def:callable-resolution}
\emph{Callable resolution} is the process of determining which
object is invoked by a call expression $f(x_1, \ldots, x_n)$.
In Python, resolution follows the Method Resolution Order (MRO)
and descriptor protocol:
\begin{enumerate}
\item Check the instance dictionary for $f$.
\item Check the class dictionary following the MRO.
\item If $f$ is a data descriptor, invoke its \texttt{\_\_get\_\_}.
\item If $f$ is in the instance dict, return it directly.
\item If $f$ is a non-data descriptor, invoke its \texttt{\_\_get\_\_}.
\end{enumerate}
\end{definition}

\begin{definition}[Binding error]
\label{def:binding-error}
A \emph{binding error} occurs when the callable resolution
produces an unexpected result, such as:
\begin{itemize}[leftmargin=2em]
\item A method is called with the wrong number of arguments.
\item A descriptor's \texttt{\_\_get\_\_} raises an exception.
\item A classmethod or staticmethod is incorrectly bound.
\item A property is called instead of accessed.
\end{itemize}
\end{definition}

The following listing shows how the $\CallAdv$ integrates with
the JuGeo API for callable analysis.

\begin{lstlisting}[style=jugeo-python, caption={Callable advisor via JuGeo API with federated verification.}, label=lst:callable-api]
from jugeo.interfaces.api import (
    JuGeoAPI,
    APIRequest,
    APIResponse,
    OperationKind,
    RequestStatus,
    CopilotAPIBridge,
    APIEventLog,
    APISession,
    APIValidator,
    APIRateLimiter,
)
from jugeo.evidence.channels import (
    CopilotChannel,
    SolverChannel,
    ChannelConfiguration,
    EvidenceChannel,
    ChannelPool,
    ChannelFederation,
)

# Full API setup for callable analysis
event_log = APIEventLog()
cop_config = ChannelConfiguration.default_for(EvidenceChannel.COPILOT)
copilot = CopilotChannel(config=cop_config)

bridge = CopilotAPIBridge(
    channel=copilot,
    event_log=event_log,
    require_corroboration=True,
)

rate_limiter = APIRateLimiter(capacity=50, refill_rate=1.0)
validator = APIValidator(max_budget=100000)

api = JuGeoAPI(
    rate_limiter=rate_limiter,
    validator=validator,
    event_log=event_log,
    copilot_bridge=bridge,
)

# Open a session for callable analysis
session = api.open_session(caller_id="callable-advisor")
assert session.is_open()

# Query for descriptor resolution analysis
descriptor_request = APIRequest(
    request_id="callable-desc-001",
    operation=OperationKind.INSPECT,
    coordinate="module.MyClass.method",
    proposition="descriptor protocol correctly applied",
    budget=15000,
    deadline=20.0,
    metadata={
        "advisor": "callable",
        "analysis": "descriptor_resolution",
    },
)

session.register_request(descriptor_request)
desc_response = api.inspect(descriptor_request)
session.complete_request(
    descriptor_request.request_id, desc_response
)

print(f"Descriptor analysis: {desc_response.status}")
print(f"Copilot contributed: {desc_response.copilot_contributed}")

# Query for binding error detection
binding_request = APIRequest(
    request_id="callable-bind-001",
    operation=OperationKind.VERIFY,
    coordinate="module.MyClass",
    proposition="all method bindings are correct",
    budget=20000,
    deadline=25.0,
    metadata={
        "advisor": "callable",
        "analysis": "binding_correctness",
    },
)

session.register_request(binding_request)
bind_response = api.verify(binding_request)
session.complete_request(binding_request.request_id, bind_response)

if bind_response.is_successful():
    print(f"All bindings correct at trust: {bind_response.trust_level}")
elif bind_response.has_residuals():
    for residual in bind_response.residuals:
        print(f"Binding residual: {residual}")

# Session history and checkpoint
history = session.recent_history()
print(f"Session requests: {len(history)}")
session.checkpoint()

# Event log export
events = event_log.export()
print(f"Events logged: {len(event_log)}")

# Bridge statistics
stats = bridge.stats()
print(f"Bridge stats: {stats}")
\end{lstlisting}

\subsection{Callable Advisory Workflow}
\label{sec:callable:workflow}

The $\CallAdv$ workflow:

\begin{enumerate}
\item \textbf{Surface enumeration.}  The LLM identifies all
  callable objects in the code under analysis.

\item \textbf{Resolution tracing.}  For each call site, the
  LLM traces the resolution chain (MRO, descriptor protocol)
  to determine the target.

\item \textbf{Error detection.}  The LLM flags potential binding
  errors, argument mismatches, and descriptor protocol violations.
\end{enumerate}

% ─────────────────────────────────────────────────────────────────────
\section{API Integration}
\label{sec:integration}

Both advisors integrate with the \JuGeo{} API through the
$\texttt{INSPECT}$ and $\texttt{VERIFY}$ operation kinds.

\begin{definition}[Advisory API protocol]
\label{def:api-protocol}
The advisory API protocol:
\begin{enumerate}
\item Open an $\texttt{APISession}$ for the advisory task.
\item Submit $\texttt{APIRequest}$s with appropriate operation kinds.
\item Receive $\texttt{APIResponse}$s with trust levels and residuals.
\item Checkpoint the session for reproducibility.
\end{enumerate}
\end{definition}

The following listing shows the composed use of both advisors
through the API.

\begin{lstlisting}[style=jugeo-python, caption={Composed import and callable advisory through the API.}, label=lst:composed]
from jugeo.interfaces.api import (
    JuGeoAPI,
    APIRequest,
    APIResponse,
    OperationKind,
    CopilotAPIBridge,
    APIEventLog,
)
from jugeo.evidence.channels import (
    CopilotChannel,
    ChannelConfiguration,
    EvidenceChannel,
    ChannelFederation,
)

# Single API instance serving both advisors
event_log = APIEventLog()
copilot = CopilotChannel(
    config=ChannelConfiguration.default_for(EvidenceChannel.COPILOT)
)
bridge = CopilotAPIBridge(
    channel=copilot, event_log=event_log, require_corroboration=True
)
api = JuGeoAPI(event_log=event_log, copilot_bridge=bridge)

# Import advisor query
import_req = APIRequest(
    request_id="composed-import-001",
    operation=OperationKind.INSPECT,
    coordinate="project.imports",
    proposition="import graph is well-formed",
    budget=10000,
    deadline=15.0,
    metadata={"advisor": "import"},
)
import_resp = api.inspect(import_req)

# Callable advisor query
callable_req = APIRequest(
    request_id="composed-callable-001",
    operation=OperationKind.INSPECT,
    coordinate="project.callables",
    proposition="all callable resolutions are correct",
    budget=10000,
    deadline=15.0,
    metadata={"advisor": "callable"},
)
callable_resp = api.inspect(callable_req)

# Federate results
federation = ChannelFederation()
all_evidence = federation.collect_all()
federation.verify_no_escalation()
composed = federation.compose_trust()

print(f"Import result: {import_resp.status}")
print(f"Callable result: {callable_resp.status}")
print(f"Composed trust: {composed}")
print(f"Total events: {len(event_log)}")
\end{lstlisting}

% ─────────────────────────────────────────────────────────────────────
\section{Formal Properties}
\label{sec:formal}

We prove five formal properties.  All proofs are in
\texttt{Paper81\_ImportCallableAdvisors.lean}.

\subsection{Import Advisory Soundness}
\label{sec:formal:import}

\begin{theorem}[Import advisory soundness]
\label{thm:import-sound}
If the $\ImportAdv$ flags an import cycle as hazardous and the
runtime channel validates the flag, then executing the import
sequence leads to an $\texttt{ImportError}$ or incomplete
module initialization.
\end{theorem}

\begin{proof}
The runtime channel validates by executing the import sequence
in an isolated environment.  If the import raises
$\texttt{ImportError}$ or produces an incompletely initialized
module (detected by attribute access on uninitialized names),
the validation confirms the hazard.  The runtime channel is
sound (it executes actual Python code), so the confirmed hazard
is genuine.
\end{proof}

\subsection{Callable Advisory Soundness}
\label{sec:formal:callable}

\begin{theorem}[Callable advisory soundness]
\label{thm:callable-sound}
If the $\CallAdv$ flags a binding error and the solver validates
the flag, then there exists an execution path where the call
produces an unexpected binding (wrong method, missing argument,
descriptor failure).
\end{theorem}

\begin{proof}
The solver constructs a formula encoding the callable resolution
and checks satisfiability of the error condition.  If satisfiable,
a model provides a concrete execution path where the binding error
manifests.  By solver soundness, this execution path is valid.
\end{proof}

\subsection{Cycle Detection Completeness}
\label{sec:formal:cycle}

\begin{definition}[Import cycle oracle]
\label{def:cycle-oracle}
The \emph{import cycle oracle} $\mathcal{O}_C$ returns all cycles
in the import graph $G$.
\end{definition}

\begin{theorem}[Cycle detection completeness (conditional)]
\label{thm:cycle-complete}
If the LLM detects all cycles in the import graph (\ie the
LLM acts as the oracle $\mathcal{O}_C$), then the $\ImportAdv$
is complete: every genuine hazardous cycle is flagged.
\end{theorem}

\begin{proof}
By the advisor pattern, every detected cycle generates an advice
item.  If the LLM detects all cycles (oracle assumption), every
cycle produces an advice item.  The classification of cycles as
hazardous or benign is a subsequent step; even if the
classification is imperfect, every cycle is at least flagged.
\end{proof}

\begin{remark}[Practical completeness gap]
\label{rem:completeness-gap}
In practice, the LLM may miss cycles that involve complex
conditional imports or dynamic import expressions.  The
completeness theorem provides an upper bound on performance.
\end{remark}

\subsection{Binding Error Monotonicity}
\label{sec:formal:binding}

\begin{theorem}[Binding error monotonicity]
\label{thm:binding-mono}
Adding a new method to a class can only increase the number of
potential binding errors: if $\mathrm{errors}(C)$ is the set of
binding errors for class $C$ and $C' = C + m$ adds method $m$:
$\mathrm{errors}(C) \subseteq \mathrm{errors}(C')$.
\end{theorem}

\begin{proof}
Adding method $m$ to class $C$ can introduce new call sites
and new descriptor resolutions, potentially creating new binding
errors.  Existing call sites may be affected if $m$ shadows an
inherited method.  However, existing errors remain errors: the
resolution that led to the error in $C$ still leads to the same
(or worse) outcome in $C'$.  Therefore
$\mathrm{errors}(C) \subseteq \mathrm{errors}(C')$.
\end{proof}

\subsection{Composition Trust Preservation}
\label{sec:formal:comp}

\begin{theorem}[Advisor composition trust]
\label{thm:comp-trust}
$\tau(\ImportAdv \AdvCompose \CallAdv) = \Tcopilot$.
\end{theorem}

\begin{proof}
Both advisors use the $\CopChan$ with ceiling $\Tcopilot$.
The composition trust is $\Tcopilot \tjoin \Tcopilot = \Tcopilot$
by idempotence of $\tjoin$.
\end{proof}

% ─────────────────────────────────────────────────────────────────────
\section{Discussion}
\label{sec:discussion}

\paragraph{Import analysis effectiveness.}
The $\ImportAdv$ is particularly effective for Python projects
with complex dependency structures.  It can identify cycles that
are difficult to spot manually, especially in large codebases
with hundreds of modules.

\paragraph{Callable analysis challenges.}
Python's dynamic nature makes callable analysis inherently
incomplete.  Dynamic dispatch, monkey-patching, and metaclass
magic can create callable surfaces that no static analysis can
fully enumerate.  The $\CallAdv$ provides best-effort analysis
guided by LLM heuristics.

\paragraph{Advisor synergy.}
The import and callable advisors complement each other: import
cycle detection can reveal modules that are partially initialized,
which in turn affects callable resolution (methods may not yet
be defined when imported cyclically).

\paragraph{Integration with IDEs.}
Both advisors can be integrated into IDE-based verification
workflows, providing real-time feedback on import and callable
issues as the programmer writes code.

% ─────────────────────────────────────────────────────────────────────
\section{Related Work}
\label{sec:related}

\paragraph{Import analysis.}
Python import analysis tools include \texttt{isort}~\cite{isort2020}
and \texttt{importlib}~\cite{python2024importlib}.  Our approach
adds LLM-guided cycle classification and trust calibration.

\paragraph{Callable analysis.}
Python callable analysis has been studied in type
checkers~\cite{mypy2024} and static analyzers.  We formalize
the callable surface and add LLM-guided error detection.

\paragraph{Advisor architectures.}
The advisor pattern in \JuGeo{} extends plugin architectures
found in IDEs~\cite{leino2010dafny} with trust calibration.

\paragraph{Graph-based program analysis.}
Import graph analysis relates to dependency analysis in build
systems~\cite{mokhov2018build} and software architecture
recovery~\cite{murphy2001software}.

% ─────────────────────────────────────────────────────────────────────
\section{Conclusion}
\label{sec:conclusion}

We have presented the $\ImportAdv$ and $\CallAdv$, two
Copilot-guided advisors for Python import graph and callable
surface analysis.  The key results are:

\begin{enumerate}
\item \textbf{Import advisory soundness} (Theorem~\ref{thm:import-sound}).
\item \textbf{Callable advisory soundness} (Theorem~\ref{thm:callable-sound}).
\item \textbf{Cycle detection completeness} (Theorem~\ref{thm:cycle-complete}).
\item \textbf{Binding error monotonicity} (Theorem~\ref{thm:binding-mono}).
\item \textbf{Composition trust preservation} (Theorem~\ref{thm:comp-trust}).
\end{enumerate}

The advisors demonstrate how LLM heuristics can guide domain-specific
program analysis within a trust-calibrated framework.

% ─────────────────────────────────────────────────────────────────────
\begin{thebibliography}{25}

\bibitem{young-jugeo}
H.~Young.
\newblock Judgment Geometry: Semantic sites, descent, and the Copilot channel.
\newblock JuGeo Paper~0, 2024.

\bibitem{young-advisors}
H.~Young.
\newblock The advisor architecture.
\newblock JuGeo Paper~90, 2024.

\bibitem{young-channels}
H.~Young.
\newblock The Copilot evidence channel.
\newblock JuGeo Paper~88, 2024.

\bibitem{young-trust}
H.~Young.
\newblock The trust algebra.
\newblock JuGeo Paper~4, 2024.

\bibitem{young-import}
H.~Young.
\newblock Import graph analysis.
\newblock JuGeo Paper~26, 2024.

\bibitem{young-callable}
H.~Young.
\newblock Callable surfaces.
\newblock JuGeo Paper~45, 2024.

\bibitem{young-heap}
H.~Young.
\newblock Heap and scope advisors.
\newblock JuGeo Paper~80, 2024.

\bibitem{chen2021codex}
M.~Chen, J.~Tworek, H.~Jun, et~al.
\newblock Evaluating large language models trained on code.
\newblock \emph{arXiv:2107.03374}, 2021.

\bibitem{isort2020}
T.~Crosley.
\newblock \texttt{isort}: A Python utility to sort imports.
\newblock \url{https://pycqa.github.io/isort/}, 2020.

\bibitem{python2024importlib}
Python Software Foundation.
\newblock \texttt{importlib} --- The implementation of import.
\newblock \emph{Python Documentation}, 2024.

\bibitem{mypy2024}
J.~Lehtosalo et~al.
\newblock \texttt{mypy}: Optional static typing for Python.
\newblock \url{https://mypy-lang.org/}, 2024.

\bibitem{leino2010dafny}
K.~R.~M.~Leino.
\newblock {Dafny}: An automatic program verifier.
\newblock In \emph{Proc.\ LPAR}, 2010.

\bibitem{mokhov2018build}
A.~Mokhov, N.~Mitchell, and S.~Peyton~Jones.
\newblock Build systems \`a la carte.
\newblock \emph{Proc.\ ACM Prog.\ Lang.}, 2(ICFP), 2018.

\bibitem{murphy2001software}
G.~C.~Murphy, D.~Notkin, and K.~J.~Sullivan.
\newblock Software reflexion models.
\newblock \emph{IEEE Trans.\ Softw.\ Eng.}, 27(4):364--380, 2001.

\bibitem{demoura2008z3}
L.~de~Moura and N.~Bj{\o}rner.
\newblock {Z3}: An efficient {SMT} solver.
\newblock In \emph{Proc.\ TACAS}, 2008.

\bibitem{de2015lean}
L.~de~Moura, S.~Kong, J.~Avigad, F.~van~Doorn, and J.~von~Raumer.
\newblock The {Lean} theorem prover.
\newblock In \emph{Proc.\ CADE}, 2015.

\bibitem{hoare1969axiomatic}
C.~A.~R.~Hoare.
\newblock An axiomatic basis for computer programming.
\newblock \emph{CACM}, 12(10):576--580, 1969.

\bibitem{ouyang2022training}
L.~Ouyang, J.~Wu, X.~Jiang, et~al.
\newblock Training language models to follow instructions.
\newblock In \emph{Proc.\ NeurIPS}, 2022.

\bibitem{dijkstra1975guarded}
E.~W.~Dijkstra.
\newblock Guarded commands.
\newblock \emph{CACM}, 18(8):453--457, 1975.

\bibitem{li2022competition}
Y.~Li, D.~Choi, J.~Chung, et~al.
\newblock Competition-level code generation.
\newblock \emph{Science}, 378(6624):1092--1097, 2022.

\end{thebibliography}

\appendix
\section{Additional Proofs}
\label{app:proofs}

\subsection{Import Graph Properties}
\label{app:import-graph}

\begin{lemma}[Cycle detection terminates]
\label{lem:cycle-term}
Cycle detection on a finite import graph $G = (V, E)$ terminates
in $O(|V| + |E|)$ time.
\end{lemma}

\begin{proof}
Standard DFS-based cycle detection visits each vertex and edge
at most once.
\end{proof}

\begin{lemma}[Acyclic import graphs are safe]
\label{lem:acyclic-safe}
If the import graph $G$ is acyclic, then all modules can be
fully initialized in topological order without circular
dependency errors.
\end{lemma}

\begin{proof}
A DAG admits a topological ordering $m_1, m_2, \ldots, m_n$
where all dependencies of $m_i$ appear before $m_i$.
Initializing modules in this order ensures each module's
dependencies are already initialized.
\end{proof}

\begin{lemma}[Star import namespace pollution bound]
\label{lem:star-bound}
If module $M$ exports $k$ names via star import, the receiving
namespace grows by at most $k$ bindings.
\end{lemma}

\begin{proof}
Each exported name either adds a new binding or overwrites an
existing one.  In the worst case, all $k$ names are new.
\end{proof}

\subsection{Callable Resolution Properties}
\label{app:callable}

\begin{lemma}[MRO is a total order on ancestor classes]
\label{lem:mro-total}
For any class $C$ in a valid Python class hierarchy, the MRO
of $C$ is a total order on the ancestors of $C$.
\end{lemma}

\begin{proof}
Python's C3 linearization produces a total order on the ancestor
classes.  If no valid linearization exists (conflicting inheritance),
Python raises a \texttt{TypeError} at class definition time.
\end{proof}

\begin{lemma}[Descriptor resolution determinism]
\label{lem:desc-determ}
For a given class hierarchy and instance state, descriptor
resolution is deterministic: the same attribute access always
resolves to the same object.
\end{lemma}

\begin{proof}
The descriptor protocol follows a fixed sequence of lookups
(data descriptor, instance dict, non-data descriptor).  Given
a fixed class hierarchy and instance state, each lookup
returns a deterministic result.
\end{proof}

\begin{lemma}[Binding error detection is decidable for finite classes]
\label{lem:binding-decidable}
For a class $C$ with finitely many methods and a finite MRO,
the set of binding errors is decidable.
\end{lemma}

\begin{proof}
Enumerate all call sites, resolve each using the MRO, and check
for argument mismatches, descriptor failures, and incorrect
bindings.  The enumeration is finite and each check is decidable.
\end{proof}

\subsection{Advisor Composition Algebra}
\label{app:algebra}

\begin{lemma}[Advisor composition forms a commutative monoid]
\label{lem:adv-monoid}
The set of advisors under $\AdvCompose$ with identity
$\mathrm{id}_{\mathrm{Adv}}$ forms a commutative monoid.
\end{lemma}

\begin{proof}
\begin{enumerate}
\item \textbf{Associativity:} $(A_1 \AdvCompose A_2) \AdvCompose A_3
  = A_1 \AdvCompose (A_2 \AdvCompose A_3)$ since union and meet
  are associative.
\item \textbf{Commutativity:} $A_1 \AdvCompose A_2 = A_2 \AdvCompose A_1$
  since union and meet are commutative.
\item \textbf{Identity:} $A \AdvCompose \mathrm{id} = A$ since
  $\mathrm{id}$ produces no advice (empty set) and has trust
  $\Tproof$ (neutral for meet).
\end{enumerate}
\end{proof}

\subsection{Import Resolution Properties}
\label{app:import-props}

\begin{definition}[Import graph]
\label{def:import-graph}
The \emph{import graph} $G_I = (M, E)$ has modules as vertices
and import edges $m \to m'$ whenever module $m$ imports module $m'$.
\end{definition}

\begin{lemma}[Import graph is a DAG for acyclic programs]
\label{lem:import-dag}
If the program has no circular imports, then $G_I$ is a directed
acyclic graph (DAG).
\end{lemma}

\begin{proof}
A circular import corresponds to a cycle in $G_I$.  If the program
has no circular imports, there are no cycles, so $G_I$ is a DAG.
\end{proof}

\begin{lemma}[Import advisor respects DAG structure]
\label{lem:import-dag-respect}
The import advisor traverses the import graph in topological order.
If $m$ imports $m'$, then $m'$ is processed before $m$.
\end{lemma}

\begin{proof}
Topological ordering of a DAG visits each vertex after all its
predecessors.  Since $m'$ is a predecessor of $m$ in the import
graph, $m'$ is processed first.  The advisor follows this ordering
when resolving imports.
\end{proof}

\begin{lemma}[Import resolution is deterministic]
\label{lem:import-determ}
For a fixed import graph and name, the resolved module is unique.
\end{lemma}

\begin{proof}
Python's import resolution follows a deterministic search path:
the local package, then \texttt{sys.path} entries in order.  The
first match is selected.  Given a fixed path, the result is unique.
\end{proof}

\begin{definition}[Callable signature space]
\label{def:callable-sig-space}
The \emph{callable signature space} $\Sigma$ is the set of all
valid Python callable signatures, including positional, keyword,
variadic, and keyword-only parameters.
\end{definition}

\begin{lemma}[Signature subsumption is a preorder]
\label{lem:sig-preorder}
The subsumption relation $\sigma_1 \preceq \sigma_2$ (meaning
$\sigma_2$ accepts all argument patterns that $\sigma_1$ accepts)
is reflexive and transitive.
\end{lemma}

\begin{proof}
\emph{Reflexive:} $\sigma \preceq \sigma$ since every signature
accepts its own argument patterns.
\emph{Transitive:} if $\sigma_1 \preceq \sigma_2$ and
$\sigma_2 \preceq \sigma_3$, then every argument pattern
accepted by $\sigma_1$ is accepted by $\sigma_2$, which
is accepted by $\sigma_3$.  So $\sigma_1 \preceq \sigma_3$.
\end{proof}

\begin{lemma}[Callable advisor type-checks suggestions]
\label{lem:callable-typecheck}
Every callable suggestion produced by the callable advisor
includes a signature that is syntactically valid Python.
\end{lemma}

\begin{proof}
The advisor parses the Copilot channel response against the
Python grammar for function signatures.  Suggestions that fail
to parse are discarded.  Therefore every emitted suggestion has
a valid signature.
\end{proof}

\begin{definition}[Call graph]
\label{def:call-graph}
The \emph{call graph} $G_C = (F, E)$ has functions as vertices
and call edges $f \to g$ whenever $f$ may call $g$ during
execution.
\end{definition}

\begin{lemma}[Call graph approximation soundness]
\label{lem:cg-sound}
The callable advisor's call graph approximation is sound:
every actual call edge is present in the approximation.
\end{lemma}

\begin{proof}
The approximation over-approximates: it includes all syntactically
possible call edges (from call sites to their possible targets
according to the type environment).  Since actual calls are a
subset of syntactically possible calls, every actual edge is
included.
\end{proof}

\begin{lemma}[Import-callable composition is well-typed]
\label{lem:import-callable-typed}
When the import advisor resolves a module and the callable advisor
analyses the resolved module's callables, the composed workflow
produces well-typed suggestions: every callable suggestion refers
to a callable that is actually defined in the resolved module.
\end{lemma}

\begin{proof}
The import advisor resolves module $m$ and provides its namespace
to the callable advisor.  The callable advisor queries the Copilot
channel with the actual namespace, so its suggestions reference
actual definitions.  Any suggestion referencing a nonexistent
callable is filtered during the parse step.
\end{proof}

\subsection{Trust and Jurisdiction Proofs}
\label{app:trust-juris}

\begin{lemma}[Import advisor jurisdiction]
\label{lem:import-juris}
The import advisor's jurisdiction covers all proposition kinds
of the form ``module $m$ exports name $n$.''
\end{lemma}

\begin{proof}
The jurisdiction is constructed by
\texttt{ChannelJurisdiction.for\_channel(EvidenceChannel.COPILOT)},
which includes all proposition kinds that the Copilot channel
can address.  Import-related propositions are within this
jurisdiction.
\end{proof}

\begin{lemma}[Callable advisor jurisdiction]
\label{lem:callable-juris}
The callable advisor's jurisdiction covers all proposition kinds
of the form ``callable $f$ has signature $\sigma$.''
\end{lemma}

\begin{proof}
Same argument as Lemma~\ref{lem:import-juris}: the Copilot
channel's jurisdiction includes callable-related propositions.
\end{proof}

\begin{lemma}[No trust escalation under composition]
\label{lem:comp-no-escal}
Composing the import advisor (trust $\Tcopilot$) with the
callable advisor (trust $\Tcopilot$) does not produce evidence
above $\Tcopilot$.
\end{lemma}

\begin{proof}
The composition trust is $\min(\Tcopilot, \Tcopilot) = \Tcopilot$.
Each advisor's output is clamped at $\Tcopilot$.  The composition
does not introduce any new evidence channel, so no escalation
occurs.
\end{proof}

\begin{proposition}[Import-callable pipeline soundness]
\label{prop:pipeline-sound}
The full import-callable pipeline (resolve imports, analyse
callables, emit suggestions) is sound: every emitted suggestion
has trust $\tleq \Tcopilot$, refers to an actual module and
callable, and has a valid signature.
\end{proposition}

\begin{proof}
Trust bound: Lemma~\ref{lem:comp-no-escal}.
Module reference: the import advisor resolves against the
actual file system (Lemma~\ref{lem:import-determ}).
Callable reference: the callable advisor analyses the resolved
namespace (Lemma~\ref{lem:import-callable-typed}).
Signature validity: Lemma~\ref{lem:callable-typecheck}.
\end{proof}

\begin{remark}[Extension to dynamic imports]
Python supports dynamic imports (\texttt{importlib.import\_module}).
The import advisor currently handles only static imports.
Extending to dynamic imports requires runtime evidence from
the runtime channel, which would increase trust to $\Truntime$
for the dynamically resolved modules.
\end{remark}

\begin{remark}[Higher-order callables]
Python callables include functions, classes (via \texttt{\_\_init\_\_}),
lambdas, and objects with \texttt{\_\_call\_\_}.  The callable
advisor handles all four cases by querying the Copilot channel
with the appropriate context.  Higher-order callables (functions
returning callables) are analysed recursively up to a configurable
depth.
\end{remark}

\subsection{Import Cycle Detection}

\begin{definition}[Import cycle]%
\label{def:importcycle}
An \emph{import cycle} is a sequence of modules
$m_1 \to m_2 \to \cdots \to m_k \to m_1$ such that
$m_i$ imports $m_{i+1}$ for each $i$.
\end{definition}

\begin{lemma}[Import advisor detects cycles]%
\label{lem:importcycledetect}
If the import graph contains a cycle of length $k$,
the import advisor detects it in $O(V + E)$ time
where $V$ and $E$ are the vertices and edges of the
import graph.
\end{lemma}
\begin{proof}
The advisor performs a depth-first traversal of the import
graph.  Back edges indicate cycles.  DFS visits each vertex
and edge once, yielding $O(V + E)$ time.
\end{proof}

\begin{lemma}[Cycle trust degradation]%
\label{lem:cycletrust}
Evidence involving a cyclic import path has trust capped at
$\Tcopilot$ regardless of the constituent channels.
\end{lemma}
\begin{proof}
Cycles introduce circular dependencies that cannot be
statically verified by the solver.  The federation falls
back to Copilot-level trust for such evidence.
\end{proof}

\subsection{Callable Arity Consistency}

\begin{lemma}[Arity-consistent suggestions]%
\label{lem:arityconsist}
The callable advisor only suggests signatures whose arity
matches the observed call sites.
\end{lemma}
\begin{proof}
The advisor collects call sites from the AST and counts
arguments at each call.  Only signatures whose parameter
count matches the observed argument counts (accounting for
default values and variadic parameters) are included in the
suggestion list.
\end{proof}

\begin{remark}[Variadic functions]
For variadic functions (\texttt{*args}, \texttt{**kwargs}),
the arity check passes for any argument count at or above
the number of positional-only parameters.  This is consistent
with Python's calling convention.
\end{remark}

\end{document}
